\begin{abstract}
Image steganography is often judged through visually convincing stego images and a small number of favourable recovery examples. Practical recoverability, however, depends on a longer and less visible chain: the transform implementation, writable coefficient coordinates, payload framing, error-control code, quantisation rule, channel definition, extraction contract, and treatment of failed decodes. When any link is underspecified, an apparently reproducible method can represent several materially different systems.

We developed an executable and auditable reconstruction framework around a contourlet-based steganography study and introduced a separate digital transport architecture designed to expose, rather than conceal, these dependencies. The architecture divides a secret image into perceptually dominant Base and Detail layers, applies unequal Reed--Solomon protection, allocates coefficients using transform energy, variance, entropy, and independently calibrated channel stability, and binds the resulting stream to its method, image pair, configuration, scrambling sequence, and interleaving pattern. Under the combined C3 configuration, the Base layer receives 65 RS(255,127) codewords rather than the mixed 32 strong and 22 weak Base codewords used by the symmetric reference. This increases the nominal Base-layer symbol-correction budget from 2,752 to 4,160 symbols, a 51.2\% increase, and raises the Base share of the protected body transport from 50.0\% to 60.2\%. The 22 formerly weak Base blocks also move from a 32-symbol to a 64-symbol correction limit.

The locked experiment comprised 16 embedding objects, 64 mandatory evaluations, and 24 conditional hard-channel checks across four fixed traceability cases. All clean cases decoded exactly. Under the evaluated JPEG, Gaussian, and salt-and-pepper conditions, however, neither the reference nor the proposed variants achieved valid end-to-end payload recovery. The experiment therefore establishes a clear code-level and architectural protection advantage for the Base layer, but not an end-to-end robustness advantage at the tested operating point. By preserving this distinction and packaging every result as a content-addressed, restart-safe research object, the study demonstrates how a negative channel result can coexist with substantial algorithmic, methodological, and reproducibility contributions.
\end{abstract}

\noindent\textbf{Keywords:} image steganography; contourlet transform; unequal error protection; adaptive embedding; Reed--Solomon coding; reproducibility; failure-aware evaluation; digital transport.

\section{Introduction}

The promise of image steganography appears simple: information should disappear into an ordinary-looking image and return intact when an authorised receiver attempts to recover it. In practice, neither half of that promise is simple. A stego image can be visually indistinguishable from its cover while the embedded payload is already fragile. Conversely, a payload can appear robust only because the method was allowed more visible distortion, a smaller message, a favourable image, or an extraction rule tuned after observing the test data.

These difficulties become sharper in transform-domain steganography. The name of a transform does not fully specify an algorithm. Filter families, directional schedules, boundary handling, coefficient ordering, subband selection, datatype conversion, and reconstruction behaviour can all alter the effective channel experienced by the payload. Two implementations may both be described as contourlet-based and still expose different coefficient sets, capacities, and failure modes.

The source study motivating this work combines geometric transformations of a secret image with contourlet-domain embedding \citep{Kumar2026}. Its direction is intuitively appealing: a directional multiscale representation can follow edges and contours more naturally than a purely separable transform, potentially offering locations where perturbations remain visually unobtrusive \citep{DoVetterli2005}. Yet several outcome-determining choices are not sufficient to identify one unique executable pipeline. Reconstructing the method therefore raised a scientific question broader than software replication: what should be claimed when the published description does not uniquely determine the experiment?

We chose not to resolve this ambiguity by silently selecting favourable assumptions. Instead, ambiguity was treated as part of the scientific object. Source-reported behaviour, engineering controls, and independently selected transform interpretations were separated. Every consequential assumption was named, versioned, hashed, and connected to the numerical evidence it produced.

That reconstruction process revealed a second problem. A transform-domain embedding rule alone is not a complete digital communication system. It does not state how the payload is framed, which information deserves priority, how errors are corrected, how corrupted output is distinguished from valid output, or how failures should be counted. We therefore introduced a digital transport path that sits above the transform and below the experimental evaluation. Its purpose is not merely to insert bits, but to make the entire journey from secret image to validated recovery explicit.

The resulting framework combines three levels of hierarchy. First, the secret is separated into Base and Detail layers according to bit significance. Second, stronger channel coding is assigned to Base. Third, in the combined adaptive configuration, Base is written before Detail into transform-coordinate groups ranked by content and calibration evidence. This creates joint semantic, coding, and spatial prioritisation.

At the same time, every method carries the same 222,360-bit protected payload and obeys the same cover--stego PSNR target of \SI{45.0}{dB} with a tolerance of \SI{0.1}{dB}. A method is therefore not permitted to appear more robust by embedding less information or by accepting more image distortion. The comparison asks a narrower question: under one fixed payload and quality budget, do adaptive allocation and unequal protection improve valid recovery?

The locked experiment produced a nuanced answer. The proposed C3 architecture is objectively stronger at the Base-layer coding level. It allocates more channel space to Base, increases its nominal correction budget by 51.2\%, doubles the correction limit of the formerly weak Base blocks, and gives Base earlier access to higher-scoring transform coordinates. Nevertheless, the evaluated channel distortions exceeded the region in which this advantage produced valid complete payload recovery. Every variant succeeded without attack, and every variant failed the strict end-to-end validity criterion under the tested medium and hard channels.

This result does not make the architecture irrelevant. It clarifies the difference between designed protection, exercised correction, partial preservation, and successful end-to-end recovery. More broadly, it demonstrates why steganography should be evaluated as an auditable digital transport rather than as a sequence of visually attractive images.

\section{Contributions}

The contribution of this work is not one isolated coefficient-selection rule. It is an integrated and auditable experimental architecture. To keep the scientific message visible, the contributions are organised into three levels.

\subsection{Primary scientific innovations}

\begin{enumerate}[leftmargin=*,label=\arabic*.]
  \item \textbf{Hierarchical Base-layer protection.} The secret is divided into perceptually dominant Base and Detail layers. C3 protects Base twice: it assigns stronger Reed--Solomon coding and places Base before Detail in transform-coordinate groups ranked by adaptive score.

  \item \textbf{Transform-bound adaptive allocation.} Eligible bands are scored using energy, variance, coefficient-magnitude entropy, and channel stability learned exclusively from separate calibration data. The calibration artifact is bound to the transform fingerprint, and robust median/MAD normalisation limits the influence of extreme bands.

  \item \textbf{A self-describing, failure-aware digital transport.} A protected header records the method, ECC layout, dimensions, layer structure, payload length, seed identifiers, payload CRC, and configuration digest. Failed Reed--Solomon or CRC validation never produces a fabricated recovered image; instead, the decoder retains layer-, stage-, and codeword-level failure evidence.

  \item \textbf{A controlled C0--C3 factorial comparison.} Adaptive allocation and unequal protection are isolated as factors A and D under a fixed 222,360-bit payload, one transform fingerprint, shared channel realisations, and a common realised cover--stego PSNR constraint.
\end{enumerate}

\subsection{Verified positive findings}

The implementation establishes several positive results that do not depend on a successful attacked decode:

\begin{itemize}[leftmargin=*]
  \item the nominal Base-layer symbol-correction budget increases from 2,752 to 4,160 symbols, a 51.2\% increase;
  \item the Base share of the protected body transport increases from 50.0\% to 60.2\%;
  \item the 22 formerly weak Base blocks move from a 32-symbol to a 64-symbol correction limit;
  \item all 16 clean C0--C3 embedding objects decode exactly under the quantised PSNR-constrained path; and
  \item C3 combines stronger Base coding with high-score-first coordinate priority.
\end{itemize}

These findings establish a verified code-level and architectural protection advantage. They are deliberately distinguished from an end-to-end robustness claim under the evaluated attacks.

\subsection{Enabling reproducibility contributions}

\begin{enumerate}[leftmargin=*,label=\alph*.]
  \item an explicit contourlet reconstruction boundary containing the 9--7 pyramid filter, PKVA directional filter, four-level schedule, P3+P4 range-coordinate pool, toolbox identity, coefficient inventory, and reconstruction gates;
  \item quality-constrained strength optimisation evaluated after inverse transformation, clipping, half-up rounding, and uint8 conversion;
  \item context-bound deterministic scrambling and interleaving derived separately for each pair, method, layer, and purpose;
  \item exact, non-overlapping coefficient maps with stable SHA-256 identities and codeword-level correction telemetry; and
  \item immutable content-addressed research objects, atomic checkpoints, cache validation, corruption quarantine, and a real SIGKILL restart gate.
\end{enumerate}

Adaptive selection, error-control coding, interleaving, and multiscale transforms have precedents in information hiding. The novelty pursued here is therefore not the isolated existence of any one ingredient. It is their precise integration under transform-bound calibration, joint coding-and-coordinate priority, fixed payload and quality constraints, strict decoding validity, and an auditable execution chain.

\section{Related Work and Scientific Positioning}

Transform-domain steganography has long used multiscale representations to balance payload capacity, visual quality, and resistance to signal processing \citep{Cox2007,Fridrich2009}. Wavelet-based methods separate coarse image structure from local detail, while the contourlet transform extends this principle through directional decomposition and improved representation of edges and contours \citep{DoVetterli2005,Daubechies1992,Mallat1999}.

The closest literature falls into three partially overlapping lines. First, contourlet steganography has previously used coefficient magnitude, contour regions, or cover selection to guide adaptive embedding \citep{SajediJamzad2010,Fakhredanesh2013}. These studies establish that contourlet structure can guide where information is hidden, but they do not define the transform-bound four-feature calibration, Base/Detail transport hierarchy, or C0--C3 factorial isolation used here. Second, unequal error protection has been applied to steganographic communication, including Turbo-code designs that assign different protection levels while maintaining an overall coding-rate constraint \citep{MaoQin2013}. Related image-transmission work has also combined wavelet importance classes with Reed--Solomon protection \citep{AlmawganiSalleh2011}. These studies establish the value of protecting more important information more strongly, but they do not jointly rank contourlet coordinates, bind calibration to a transform fingerprint, and expose codeword-level recovery evidence under one fixed steganographic payload. Third, the source study combines AP/GP/HP transformations with LPDFB contourlet embedding and reports robustness under several signal-processing and geometric attacks \citep{Kumar2026}. It does not provide the layered digital transport, prospective ablation, or object-level provenance introduced here.

\begin{table}[htbp]
\centering
\caption{Closest prior-work lines and the specific novelty boundary of the present study. The table does not claim that adaptive embedding, UEP, Reed--Solomon coding, or contourlet analysis is individually unprecedented.}
\label{tab:novelty}
\small
\begin{tabularx}{\textwidth}{>{\raggedright\arraybackslash}p{0.17\textwidth} >{\raggedright\arraybackslash}X >{\raggedright\arraybackslash}X >{\raggedright\arraybackslash}X}
\toprule
Prior-work line & Established contribution & Missing relative to the present question & Present extension \\
\midrule
Kumar et al. \citep{Kumar2026}
& AP/GP/HP transformation of the secret and LPDFB contourlet embedding; reported imperceptibility and attack results.
& Outcome-determining transform and transport details are not fully harmonised for an author-equivalent comparison; no C0--C3 ablation or layered ECC transport is described.
& Explicit independent PDFB interpretation, fixed payload and PSNR contract, layered transport, failure-aware validity, and immutable evidence chain. \\
\addlinespace
Adaptive contourlet steganography \citep{SajediJamzad2010,Fakhredanesh2013}
& Contourlet magnitude or contour structure guides cover/region selection and adaptive embedding.
& No separate attack-calibration artifact bound to the transform fingerprint; no joint Base/Detail coding hierarchy or factorial A/D isolation.
& Four-feature transform-bound scoring, robust normalisation, calibration-gated execution, and high-score-first Base placement. \\
\addlinespace
UEP for steganographic communication \citep{MaoQin2013}
& Different protection levels are assigned to information classes under a coding-rate constraint.
& Does not couple UEP to contourlet-coordinate ranking, fixed digital stego quality, or layer/codeword failure telemetry.
& Reed--Solomon Base/Detail UEP combined with adaptive coordinate priority and strict CRC-validated recovery. \\
\addlinespace
RS-based UEP image transmission \citep{AlmawganiSalleh2011}
& Important wavelet information receives stronger RS-based protection in an error-prone channel.
& Image transmission rather than hidden-payload transport; no steganographic coefficient map, self-describing stream, or cover--stego quality constraint.
& Integrates importance-aware RS protection into a fixed-length contourlet steganographic transport with exact coefficient identities. \\
\bottomrule
\end{tabularx}
\end{table}

Table~\ref{tab:novelty} identifies the novelty at the level most likely to matter in peer review: not ``the first use'' of familiar components, but a previously unreported experimental combination and evidence contract. The proposed path defines C0--C3 as a prospective \(2\times2\) design: adaptive allocation is factor A, unequal protection is factor D, and C3 combines both. This makes the internal comparison mechanistic rather than promotional.

Direct numerical superiority over the published source article remains outside the present evidence boundary. Transform parameters, secret sizes, payload contracts, image pairing, extraction assumptions, quantisation rules, and some attack details are not fully harmonised. The paper therefore distinguishes three claims that are often blurred: an executable independent reconstruction, an internally controlled new transport architecture, and an author-equivalent reproduction. Only the first two are supported here.
